\section{Conclusions}
\label{sec:conclusions}

We introduced an open-weight, open-source framework for translating natural-language strategic requirements into well-formed ATL/ATL$^\ast$ specifications. We provided an expert-curated dataset explicitly designed for the NL-to-ATL/ATL$^\ast$ task and covering both simple requirements and five ambiguity families: right dislocation, left dislocation, VP ellipsis, Right Node Raising, and Quantifier Scope Ambiguity. The dataset is approximately balanced with respect to ATL/ATL$^\ast$ temporal operators and provides a reusable benchmark for training and evaluating models for strategic specification synthesis. We also presented a hybrid framework running both cloud APIs and local open-weight models with LoRA fine-tuning; and a multi-faceted evaluation protocol coupling normalized exact match with a human-validated LLM-as-a-judge. Our experiments show that in-domain fine-tuning of small, locally deployable open-weight models ($3$-$7$B) translates strategic requirements as accurately as strong few-shot proprietary baselines under the most human-aligned judge; that exact match alone understates quality, since faithful translations often differ syntactically from the reference; and that the automatic judges must themselves be validated, since judge reliability runs opposite to generator capability, the open-weight Llama-3.3-70B is the most human-aligned judge, while the strongest, most recent proprietary models are the least reliable, over-rejecting faithful paraphrases even of their own generations. Finally, the integration with \textsc{VITAMIN} shows how natural-language strategic requirements can be connected to existing model-checking workflows.

%\paragraph{Limitations.}
%Several limitations qualify these results. The benchmark, although expert-validated, is curated rather than drawn from deployed industrial requirements, and is English-only; transfer to in-the-wild specifications and other languages remains to be demonstrated. We target ATL/ATL$^\ast$, so requirements involving knowledge, beliefs, or the more expressive constructs of Strategy Logic lie outside the present scope. Semantic accuracy is measured by an LLM judge whose agreement with human experts (moderate even for the best-aligned judge) is imperfect, so the absolute numbers are best read as calibrated estimates rather than exact rates; we mitigate this with multiple independent judges and a human-overridden ranking. Our reported confidence intervals, moreover, reflect training-seed variation at a single canonical split; with only $218$ test items, test-set resampling adds uncertainty of order $\pm0.06$ that these intervals do not capture, so they are best read as a lower bound, a stratified $k$-fold evaluation, supported by our framework, would tighten this estimate. Finally, equivalence is assessed by strategic-temporal intent rather than by formal model checking, and the reported cross-environment latencies are only indicative across the API/local boundary.

\paragraph{Future Work.}
This seminal study provides a springboard for several future research directions. The main avenue concerns extending the dataset by releasing it publicly to the formal methods community for their contribution, as well as broadening the framework to more expressive strategic formalisms, such as Strategy Logic and its fragments~\cite{DBLP:journals/tocl/MogaveroMPV14,DBLP:conf/aaai/CermakLM15,DBLP:conf/ijcai/BelardinelliJKM19}, or epistemic variants of ATL~\cite{DBLP:journals/fuin/JamrogaH04,aminof2016prompt,jamroga2017reasoning}. This would enable translations involving knowledge, beliefs, and information flow in multi-agent systems. Moreover, the released dataset opens the door to training domain-specialized models potentially reducing reliance on post-hoc semantic validation.

%Future work will focus on interactive disambiguation, richer strategic formalisms, and deeper coupling between translation models and verification feedback.
